% Source PDF: source/普通高考/2010/2010辽宁文.pdf, page 1, question 5.
% pdfimages -list reports no embedded bitmap; the flowchart is vector/text artwork.
% Grid evidence: tmp/review_flowchart_2010_liaoning_liberal/grid/page-1-grid100.png
% (100px coarse + 50px fine) and q5-block-grid50.png (50px coarse + 25px fine).
% Comparison crop: tmp/review_flowchart_2010_liaoning_liberal/crop/q5_original_final.png,
% taken from the ungridded page render with a safety border.
% Nodes: 开始; 输入 n,m; k=1,p=1; p=p(n-m+k); k<m?; k=k+1; 输出 p; 结束.
% Directed edges: 开始→输入→初始化→乘积更新→判定;
% 判定(是)→k更新→上回线左向箭头→乘积更新; 判定(否)→输出→结束.
% The loop return is independent of the output branch; all node text is centered
% with local zero paragraph indentation.

\begin{tikzpicture}[
  x=1cm,y=1cm,baseline,
  flowarrow/.style={-{Latex[length=2.1mm,width=1.5mm]},line width=.45pt},
  nodebox/.style={draw,line width=.45pt,inner sep=2pt,minimum height=.68cm,
    anchor=center,align=center,execute at begin node={\setlength{\parindent}{0pt}\noindent}},
  term/.style={nodebox,rounded corners=2.5pt,minimum width=1.30cm,text width=1.30cm},
  io/.style={nodebox,shape=trapezium,trapezium left angle=70,trapezium right angle=110,
    minimum width=2.30cm,text width=1.90cm,minimum height=.72cm},
  decision/.style={nodebox,diamond,aspect=2.75,minimum width=3.55cm,minimum height=1.02cm,inner sep=1pt}
]
  \node[term] (start) at (0,10.00) {开始};
  \node[io,minimum width=2.45cm,text width=2.05cm] (input) at (0,8.75) {输入 $n,m$};
  \node[nodebox,text width=3.30cm,minimum width=3.30cm] (init) at (0,7.45)
    {$k=1,p=1$};
  \node[nodebox,text width=4.20cm,minimum width=4.20cm] (product) at (0,6.20)
    {$p=p(n-m+k)$};
  \node[decision] (test) at (0,4.55) {$k<m$};
  \node[nodebox,text width=2.65cm,minimum width=2.65cm] (increment) at (3.60,5.60)
    {$k=k+1$};
  \node[io,minimum width=2.20cm] (output) at (0,3.00) {输出 $p$};
  \node[term] (finish) at (0,1.55) {结束};

  \coordinate (loopjoin) at (0,6.80);
  \coordinate (branchright) at (3.60,4.55);
  % Keep the return leg strictly orthogonal: rise above the increment box,
  % then run left to the shared point on the main vertical stem.
  \coordinate (righttop) at (3.60,6.80);

  % Main stem and unique downward arrows.
  \draw[flowarrow] (start.south) -- (input.north);
  \draw[flowarrow] (input.south) -- (init.north);
  \draw[flowarrow] (init.south) -- (product.north);
  \draw[flowarrow] (product.south) -- (test.north);

  % Yes branch: enter the increment node from below, then return above the
  % product node and point left into the shared stem.
  \draw (test.east) -- node[pos=.65,above=1pt] {是} (branchright);
  \draw[flowarrow] (branchright) -- (increment.south);
  \draw[flowarrow] (increment.north) -- (righttop) -- (loopjoin);

  % No branch exits to output and termination, independent of the loop.
  \draw[flowarrow] (test.south) -- node[right=3pt,pos=.45] {否} (output.north);
  \draw[flowarrow] (output.south) -- (finish.north);

  \path[use as bounding box] (-2.55,1.15) rectangle (5.05,10.45);
\end{tikzpicture}
